\section{Related Work}
\label{sec:related}

\subsection{Knowledge Tracing Models and Benchmarking}
Knowledge Tracing models dynamic learner proficiency by leveraging historical interaction sequences to forecast future responses \citep{abdelrahman2023knowledge}. Classical approaches, such as Item Response Theory (IRT) \citep{rasch1960probabilistic,embretson2000item} and Bayesian Knowledge Tracing (BKT) \citep{corbett1994knowledge}, model static ability and latent mastery states, respectively. Deep Knowledge Tracing (DKT) \citep{piech2015deep} shifted the field toward recurrent neural networks, followed by memory-augmented architectures such as DKVMN \citep{zhang2017dynamic} and attention-based models such as AKT \citep{ghosh2020context} and SimpleKT \citep{liu2023simplekt}. Benchmarking toolkits such as pyKT \citep{liu2022pykt} have substantially improved the standardization of preprocessing and model evaluation, and have documented how easily evaluation choices can distort reported performance. Our work is complementary to this benchmarking line: rather than standardizing model training and comparison, we standardize diagnostic reporting---how predictive and probabilistic quality should be broken down, audited, and interpreted once predictions have been produced. We therefore use representative baselines from this literature to instantiate the protocol, not to compete with them.

\subsection{Graph-Augmented Knowledge Tracing}
Contemporary architectures frequently leverage graph neural networks to capture structural dependencies---such as prerequisites, semantic similarities, or co-occurrences---between knowledge components (KCs) \citep{abdelrahman2023knowledge}. GKT \citep{nakagawa2019graph} and GIKT \citep{yang2021gikt} construct KC--item or KC--KC dependency graphs, while contrastive and self-supervised approaches such as CL4KT \citep{lee2022cl4kt} learn from structurally augmented views. These architectures capture rich relational signal, but they also enlarge the evaluation attack surface: graph construction, augmentation, and self-supervised objectives introduce additional channels through which information can leak across splits, and their intended gains are typically claimed precisely on low-frequency concepts, where evaluation is least stable. Reliable stratum-level and leakage-audited evaluation is therefore a prerequisite for credible claims in this family. Our protocol deliberately instantiates the diagnostics on standard sequential baselines to establish that foundation; extending the same audited reporting to graph-based architectures is a natural next step.

\subsection{Sparse and Cold-start Knowledge Components}
Educational datasets exhibit long-tail distributions in which a small set of dense KCs dominates the interaction volume while many KCs carry sparse training evidence \citep{choi2020ednet,liu2023xes3g5m}. Sparse and cold-start conditions in KT should be distinguished from the cold-start scenarios typically studied in recommender systems: recommender systems usually define cold-start at the user or item level (new users or items)---a challenge also studied in KT as the new student problem \citep{bhattacharjee2025coldstart}---whereas the variant of the problem addressed here concerns concepts with insufficient interaction history in the training data \citep{pelanek2015memory,choffin2019das3h}. Our protocol operationalizes this distinction through train-only KC-frequency stratification that separates strict cold-start concepts ($f_{\mathrm{train}}(c) = 0$) from very sparse, sparse, medium, and dense strata, so that concepts with no training evidence are never conflated with concepts that have limited but nonzero evidence.

The term sparse has also been used in a different, model-internal sense. sparseKT \citep{huang2023sparsekt} introduces $k$-sparse attention that selects the most relevant historical interactions to improve robustness and generalization. This is orthogonal to our usage: sparseKT sparsifies attention over interaction histories, whereas our protocol diagnoses low-frequency and zero-frequency KCs via train-only counts. Sparse attention and sparse-concept evaluation should not be conflated.

Recent work addresses cold-start KT from complementary angles. csKT \citep{bai2025cskt} targets short interaction sequences via kernel bias and cone attention; CLST \citep{jung2025clst} reformulates KT with natural-language descriptions of exercises and KCs; and further LLM-based approaches include attribute-aware KT \citep{guo2024mitigating} and option-weighted KT \citep{kim2024beyond}. These methods target learner-level, sequence-level, system-level, or semantic cold-start, whereas our protocol defines and audits cold-start at the concept level through train-only frequency. Because this paper proposes a diagnostic protocol rather than a new architecture or attention mechanism, we do not compare against these model-level methods; instead, the proposed diagnostics can be applied to them---sparseKT-like and cold-start KT models alike---in future evaluations.

\subsection{Calibration of Probabilistic Predictions}
Calibration is essential for KT because predicted probabilities of correctness can directly drive instructional interventions. In student modeling, Pel\'{a}nek's evaluation framework is the central prior reference for probability-level assessment using the Brier score and calibration curves \citep{pelanek2015metrics}, establishing that student models should be evaluated beyond binary accuracy. In the broader machine-learning literature, deep networks can achieve high discrimination while remaining poorly calibrated and overconfident \citep{guo2017calibration}; calibration error is commonly quantified via the Expected Calibration Error (ECE) \citep{naeini2015obtaining}, and the Brier score \citep{brier1950verification} admits a decomposition into uncertainty, reliability, and resolution \citep{degroot1983comparison} that separates intrinsic task difficulty from model miscalibration.

Our protocol builds directly on this tradition, and differs from it in how the calibration evidence is organized: prior calibration-oriented student-model evaluation is typically reported at the population level, rather than around train-only KC-frequency strata, an explicit concept-level strict cold-start group, and sample-size-aware bucket diagnostics with reliability flags. To the best of our knowledge, this combination---per-stratum ECE, Brier decomposition, and reliability diagrams, gated by sample-size-aware reporting and a leakage audit---has not been packaged as a reusable, dataset-agnostic protocol for KT evaluation. Complementary to our focus on evaluation diagnostics, recent work has begun to explore methodological corrections for per-item bias and stranded discrimination \citep{yan2026recovering}; our protocol provides the diagnostic foundation to assess whether such corrections effectively improve reliability across varying frequency strata.

\subsection{Reproducibility, Leakage Control, and Statistical Testing}
As sequence modeling architectures grow increasingly complex, ensuring rigorous reproducibility and strict leakage control has become critical for reliable evaluation \citep{kapoor2023leakage, gervet2020deep, kim2026ensuring}. The field faces a leakage and reproducibility crisis, and KT is particularly susceptible to silent evaluation errors due to its complex sequence structure. To mitigate this, our protocol's L1--L8 audit defines named channels tied to verifiable artifacts. For comparing competitive models, DeLong's test provides a standard approach for correlated ROC curves \citep{delong1988comparing}. Statistical testing is used only where prediction-level outputs permit paired comparison; otherwise, we report descriptive diagnostics with sample-size-aware interpretation.
